\documentclass[11pt,a4paper]{article}

\usepackage[T1]{fontenc}
\usepackage[utf8]{inputenc}
\usepackage{lmodern}
\usepackage{amssymb,amsmath,amsthm,mathtools}
\usepackage[margin=2.5cm]{geometry}
\usepackage{parskip}
\usepackage[hidelinks]{hyperref}
\usepackage{microtype}

% Theorem environments (arXiv / journal style)
\theoremstyle{plain}
\newtheorem{theorem}{Theorem}[section]
\newtheorem{proposition}[theorem]{Proposition}
\newtheorem{lemma}[theorem]{Lemma}
\newtheorem{corollary}[theorem]{Corollary}

\theoremstyle{definition}
\newtheorem{definition}[theorem]{Definition}
\newtheorem{example}[theorem]{Example}
\newtheorem{remark}[theorem]{Remark}

% Notation
\newcommand{\Type}{\mathbf{Type}}
\newcommand{\id}{\mathrm{id}}
\newcommand{\Bom}{\mathsf{Bom}}
\newcommand{\bom}{\mathsf{bom}}
\newcommand{\Occ}{\mathsf{Occ}}
\newcommand{\site}{\mathsf{site}}
\newcommand{\op}{\mathrm{op}}

\title{Sourcing as a Functor of Directed Containers:\\
  Routed Bills of Materials}
\author{Extracted from a machine-checked development}
\date{}

\begin{document}
\maketitle

\begin{abstract}
We formalise the sourcing assignment of a bill of materials (BOM) as a
covariant functor of directed containers. The domain is the directed
container of \emph{routed} BOMs: finite product-labelled rooted trees whose
edges carry inbound shipping routes. The codomain is the shipping network
read in the opposite direction---positions at a facility are inbound paths
ending there. On shapes, the functor reads off the facility producing the
root product; on positions, an occurrence is sent to the inbound route
composed leg-by-leg through intermediate producers. Functoriality is
\emph{consolidation}: the empty occurrence is sourced by the empty route,
and nested occurrences concatenate their routes. The proofs of all stated
results have been machine-checked; we record only the mathematical content.
\end{abstract}

\tableofcontents
\bigskip

%======================================================================
\section{Introduction}
%======================================================================

A bill of materials is a tree of products: the finished good at the root,
direct components as children, raw materials at the leaves. In a physical
supply chain each product is made at some facility, and each component
travels an inbound route to the point of use of its parent. The assignment
of sites and inbound routes is the \emph{sourcing} data of the BOM.

Abstractly, both the family of BOMs and the shipping network may be
presented as directed containers. Sourcing then becomes a morphism of that
structure---more precisely, a covariant functor of directed containers
from the BOM container to the opposite shipping network. The content of
this note is to make that functor precise, and to record the
formalisation choice that makes it total.

\paragraph{The formalisation hypothesis.}
If the shapes of the BOM container were merely product-labelled trees,
with no route data, a functor to the opposite network would have to
produce an inbound path for every tree edge---including perverse trees
whose edges run against the shipping graph. No such path need exist, so
the data cannot extend to a total functor on all trees.

The remedy adopted here is to place the routing choices on the shapes
themselves. A \emph{routed BOM} is a product-labelled rooted tree indexed
by the facility producing its root, in which every edge already carries
an inbound route. On this domain the sourcing functor is canonical: no
side data, and the functoriality equations hold definitionally (or by a
single path induction for consolidation). The paper's informal data---a
site per product and one leg per BOM edge---is exactly what builds a
concrete routed BOM.

\paragraph{What this note claims.}
In ordinary mathematical language: routed BOMs form a directed container
whose positions are occurrences (descending paths from the root). The
shipping digraph~$G$, read oppositely, forms a directed container whose
positions at a vertex are inbound routes. Sourcing is the functor that
sends a routed BOM to its root facility and an occurrence to the
composite inbound route of that component. Functoriality is consolidation
of multi-level routes. The UAE steel rebar example of the companion paper
is a concrete instance, with the three-leg ore-to-rolling-mill route
checked by computation.

We work throughout in a single universe of types, writing $\Type$ for that
universe. Proofs are indicated by the phrase ``by lean,'' meaning they
have been checked by a proof assistant rather than written out in prose.

%======================================================================
\section{Background: digraphs, paths, and opposite networks}
%======================================================================

\begin{definition}[Digraph]\label{def:digraph}
A \textbf{digraph} $G$ consists of a type $V_G$ of vertices, a type $E_G$
of edges, and source and target maps
\[
  \mathsf{src}_G,\;\mathsf{tgt}_G \;:\; E_G \;\longrightarrow\; V_G.
\]
\end{definition}

\begin{definition}[Path]\label{def:path}
A \textbf{path} in~$G$ from $u$ to $v$ is a finite walk of edges with
matching endpoints. Paths form a category: there is an empty path
$\mathsf{nil}_u$ at each vertex, and concatenation
\[
  \mathsf{append} \;:\;
  \mathsf{Path}_G(u,v) \times \mathsf{Path}_G(v,w)
  \;\longrightarrow\;
  \mathsf{Path}_G(u,w)
\]
is associative, with $\mathsf{nil}$ as two-sided unit.
\end{definition}

\begin{definition}[Opposite digraph]\label{def:op}
The \textbf{opposite digraph} $G^\op$ has the same vertices and edges as
$G$, with source and target swapped:
\[
  \mathsf{src}_{G^\op} = \mathsf{tgt}_G,\qquad
  \mathsf{tgt}_{G^\op} = \mathsf{src}_G.
\]
Consequently a path of $G^\op$ from $v$ to $u$ is an \emph{inbound route}
of~$G$ ending at~$v$ and originating at~$u$.
\end{definition}

\begin{remark}
The opposite network is taken on the graph, not as a general dual of
directed containers: one forms the directed container associated to
$G^\op$ in the usual way (shapes are vertices; positions at~$v$ are paths
of $G^\op$ starting at~$v$). Positions at a facility are therefore inbound
routes to that facility, read backwards from the point of use.
\end{remark}

%======================================================================
\section{Routed bills of materials}
%======================================================================

Fix a type $P$ of products and a digraph $G$ of shipping facilities and
legs.

\begin{definition}[Routed BOM]\label{def:bom}
A \textbf{routed BOM} over $P$ and $G$, indexed by a facility
$v \in V_G$, is a finite tree
\[
  t \;=\; \mathsf{node}\bigl(
    \mathit{prod},\; v,\; n,\;
    (v_i)_{i < n},\;
    (\rho_i)_{i < n},\;
    (t_i)_{i < n}
  \bigr)
\]
where
\begin{enumerate}
  \item $\mathit{prod} \in P$ is the product label of the root,
  \item $n \in \mathbb{N}$ is the number of direct components,
  \item $v_i \in V_G$ is the facility producing the $i$-th component,
  \item $\rho_i$ is a path of $G^\op$ from $v$ to $v_i$ (the inbound route
        of that component into the root facility),
  \item $t_i$ is a routed BOM at $v_i$ (the component's own BOM).
\end{enumerate}
Raw materials are nodes with $n = 0$. We write $\Bom(P,G)(v)$ for the
type of routed BOMs at~$v$.
\end{definition}

Thus a routed BOM is a product-labelled rooted tree whose every edge
already carries its sourcing route. The routes live on the shapes so that
the sourcing functor below is total.

\begin{definition}[Occurrence]\label{def:occ}
An \textbf{occurrence} in a routed BOM $t$ is a descending path from the
root, including the empty path:
\begin{itemize}
  \item $\mathsf{here}$ is the empty occurrence (the root itself);
  \item if $i < n$ and $p$ is an occurrence in the $i$-th child, then
        $\mathsf{inside}(i,p)$ is an occurrence in the parent.
\end{itemize}
We write $\Occ(t)$ for the type of occurrences in~$t$. Labels, sites, and
routes play no role in the navigation.
\end{definition}

\begin{definition}[Sub-BOM and shift]\label{def:sub-shift}
\begin{enumerate}
  \item The \textbf{sub-BOM} at an occurrence $p \in \Occ(t)$ is the
        subtree $t.\mathsf{sub}\,p$, together with the facility
        producing it.
  \item The \textbf{shift} (occurrence concatenation)
        \[
          t.\mathsf{shift} \;:\;
          (p : \Occ(t)) \;\longrightarrow\;
          \Occ(t.\mathsf{sub}\,p) \;\longrightarrow\;
          \Occ(t)
        \]
        views an occurrence inside a component as an occurrence of the
        finished product.
\end{enumerate}
\end{definition}

\begin{proposition}[Navigation laws]\label{prop:nav}
For every routed BOM $t$ and occurrences $p$, $q$, $r$ of the appropriate
types:
\begin{enumerate}
  \item[\textup{(D3)}] $t.\mathsf{shift}\,p\,\mathsf{here} = p$;
  \item[\textup{(D4)}] $t.\mathsf{sub}\,(t.\mathsf{shift}\,p\,q)
        = (t.\mathsf{sub}\,p).\mathsf{sub}\,q$;
  \item[\textup{(D5)}] shift is associative along the transport of
        \textup{(D4)}.
\end{enumerate}
\end{proposition}

\begin{proof}
By lean.
\end{proof}

\begin{remark}
Law~(D4) is the slogan of multi-level BOM explosion: sub-BOMs nest along
the shift, so multi-level costing is coherent. Law~(D5) is the
associativity of that nesting.
\end{remark}

%======================================================================
\section{The directed container of routed BOMs}
%======================================================================

\begin{definition}[BOM container]\label{def:bom-dc}
The \textbf{directed container of routed BOMs} $\bom(P,G)$ has
\begin{enumerate}
  \item shapes $(v,t)$ with $v \in V_G$ and $t \in \Bom(P,G)(v)$;
  \item positions of $(v,t)$ equal to $\Occ(t)$;
  \item root position $\mathsf{here}$;
  \item sub-shape and shift as in Definition~\ref{def:sub-shift}.
\end{enumerate}
The five directed-container axioms hold by the navigation laws of
Proposition~\ref{prop:nav} (together with the definitional equalities for
root).
\end{definition}

\begin{remark}
Labels, sites, and routes ride along on the shapes without disturbing the
navigation laws. The combinatorial skeleton is that of rooted trees; the
supply-chain decoration is passive data.
\end{remark}

%======================================================================
\section{The sourcing functor}
%======================================================================

\begin{definition}[Inbound route of an occurrence]\label{def:routeOf}
Let $t$ be a routed BOM at facility~$v$. The \textbf{inbound route} of an
occurrence $p \in \Occ(t)$ is the path of $G^\op$ from $v$ to the facility
producing the sub-BOM at~$p$, defined by induction on~$p$:
\begin{align}
  t.\mathsf{routeOf}(\mathsf{here})
  &\;=\;
  \mathsf{nil}_v, \\
  t.\mathsf{routeOf}\bigl(\mathsf{inside}(i,p)\bigr)
  &\;=\;
  \rho_i \mathbin{+\mkern-6mu+}
  \bigl((t_i).\mathsf{routeOf}\,p\bigr),
\end{align}
where $\rho_i$ is the inbound leg stored on the $i$-th edge and
$+\mkern-6mu+$ denotes path concatenation. Thus the empty occurrence is
sourced by the empty route, and a nested occurrence composes leg-by-leg
through intermediate producers.
\end{definition}

\begin{theorem}[Consolidation]\label{thm:consolidation}
For every routed BOM $t$ and occurrences $p$, $q$ of the appropriate
types, the route of a concatenated occurrence is the concatenation of the
routes:
\begin{equation}
  \bigl\langle
    (t.\mathsf{sub}\,(t.\mathsf{shift}\,p\,q)).1,\;
    t.\mathsf{routeOf}\,(t.\mathsf{shift}\,p\,q)
  \bigr\rangle
  \;=\;
  \bigl\langle
    ((t.\mathsf{sub}\,p).\mathsf{sub}\,q).1,\;
    (t.\mathsf{routeOf}\,p)
    \mathbin{+\mkern-6mu+}
    ((t.\mathsf{sub}\,p).\mathsf{routeOf}\,q)
  \bigr\rangle,
\end{equation}
as positions of the opposite network (endpoint paired with route).
\end{theorem}

\begin{proof}
By lean.
\end{proof}

\begin{remark}
The two endpoints agree by~(D4), but not definitionally; packaging them
with their routes as dependent pairs lets the equality be stated without
transport. Consolidation is the position-level form of multi-level
sourcing: ore reaches the rolling mill only via the reduction plant and
the melt shop.
\end{remark}

\begin{definition}[Sourcing functor]\label{def:sourcing}
The \textbf{sourcing functor}
\[
  \site_{P,G}
  \;:\;
  \bom(P,G)
  \;\longrightarrow\;
  \mathsf{network}(G^\op)
\]
is the covariant functor of directed containers given by
\begin{enumerate}
  \item on shapes: $(v,t) \mapsto v$ (the facility producing the root
        product);
  \item on positions: $p \mapsto
        \bigl\langle
          (t.\mathsf{sub}\,p).1,\;
          t.\mathsf{routeOf}\,p
        \bigr\rangle$
        (the inbound route of the occurrence, ending at the component's
        own producing facility).
\end{enumerate}
\end{definition}

\begin{theorem}[Functoriality]\label{thm:functor}
The maps of Definition~\ref{def:sourcing} satisfy the three axioms of a
covariant functor of directed containers:
\begin{enumerate}
  \item[\textup{(F0)}] sub-shapes are preserved (the network's sub is
        endpoint projection);
  \item[\textup{(F1)}] the root occurrence is sent to the empty route;
  \item[\textup{(F2)}] shift is preserved --- this is consolidation
        (Theorem~\ref{thm:consolidation}).
\end{enumerate}
\end{theorem}

\begin{proof}
By lean.
\end{proof}

\begin{remark}
Axioms~(F0) and~(F1) hold definitionally. Axiom~(F2) is the single
non-trivial content; its transport vanishes because~(F0) is
definitional. There is no side data: once a BOM is routed, sourcing is
canonical.
\end{remark}

%======================================================================
\section{Example: UAE steel rebar}
%======================================================================

We record the concrete instance that motivates the construction: a
gas-based mini-mill making rebar for export.

\begin{example}[Steel network and BOM]\label{ex:steel}
The shipping graph is the directed path of four facilities
\[
  \text{ore berth}
  \;\longrightarrow\;
  \text{reduction plant}
  \;\longrightarrow\;
  \text{melt shop}
  \;\longrightarrow\;
  \text{rolling mill},
\]
with three legs $\ell_{\mathrm{ob}\to\mathrm{rp}}$,
$\ell_{\mathrm{rp}\to\mathrm{ms}}$,
$\ell_{\mathrm{ms}\to\mathrm{rm}}$.

The products are ore pellets, DRI, scrap, billet, and rebar. The routed
BOM at the rolling mill is:
\begin{itemize}
  \item rebar at the rolling mill, with billets inbound along
        $\ell_{\mathrm{ms}\to\mathrm{rm}}$;
  \item billets at the melt shop, with DRI inbound along
        $\ell_{\mathrm{rp}\to\mathrm{ms}}$ and scrap already on site
        (empty inbound route);
  \item DRI at the reduction plant, with ore pellets inbound along
        $\ell_{\mathrm{ob}\to\mathrm{rp}}$;
  \item ore pellets at the ore berth, and scrap at the melt shop, as
        leaves.
\end{itemize}
\end{example}

\begin{proposition}[Sourcing on the steel instance]\label{prop:steel}
Let $t$ be the rebar BOM of Example~\ref{ex:steel}, viewed as a shape of
$\bom(P,G)$.
\begin{enumerate}
  \item On shapes, $\site(t)$ is the rolling mill.
  \item The ore-pellet occurrence is sourced by the three-leg inbound
        route
        \[
          \text{rolling mill}
          \;\leftarrow\;
          \text{melt shop}
          \;\leftarrow\;
          \text{reduction plant}
          \;\leftarrow\;
          \text{ore berth}.
        \]
  \item The scrap occurrence is sourced by the one-leg route ending at
        the melt shop.
  \item Consolidation holds on the instance: sourcing the concatenated
        occurrence ``billet then ore'' equals shifting the billet's route
        by the ore's route in the opposite network.
\end{enumerate}
\end{proposition}

\begin{proof}
By lean.
\end{proof}

%======================================================================
\section{Discussion}
%======================================================================

\paragraph{What has been established.}
Routed BOMs form a directed container; the opposite shipping network is a
directed container; sourcing is a covariant functor between them whose
position action is multi-level route composition. The UAE steel example
confirms the equations by direct computation.

\paragraph{Why routes live on shapes.}
The informal sourcing data of a companion development---a site per
product and one inbound leg per BOM edge---cannot extend to a total
functor on \emph{all} product-labelled trees, because some trees demand
routes that do not exist in~$G$. Equipping each shape with its routing
choices restores totality and yields a canonical functor. An alternative
(restricting to trees compatible with a fixed component relation) makes
shape equalities propositional and reintroduces transports throughout.

\paragraph{Scope.}
The development treats a single universe of types and does not address
costing, lot sizes, or capacity constraints. Directed-container structure
records only the combinatorial skeleton of navigation and sourcing; the
economic content rides along as passive data on shapes and positions.

\paragraph{Machine checking.}
Every definition and theorem statement above corresponds to a declaration
in a machine-checked formalisation. The phrase ``by lean'' indicates that
the proof has been verified by a proof assistant; we have not reproduced
the formal proof scripts here.

%======================================================================
\section*{Acknowledgements}
%======================================================================

The mathematical content follows the theory of directed containers and
the supply-chain reading developed in the companion notes on supply chains
as directed containers. The presentation here is a prose rendering of a
formal development intended for human review of the theorem statements.

\begin{thebibliography}{9}

\bibitem{ahman2016}
D.~Ahman and T.~Uustalu.
\newblock Directed containers as categories.
\newblock In \emph{MSFP}, volume~207 of \emph{EPTCS}, pages 89--98, 2016.

\bibitem{abbott2005}
M.~Abbott, T.~Altenkirch, and N.~Ghani.
\newblock Containers: constructing strictly positive types.
\newblock \emph{Theoretical Computer Science}, 342(1):3--27, 2005.

\end{thebibliography}

\end{document}
